
\subsubsection{Key Findings}

Model rankings are highly stable between ImageNet and ImageNet-V2 ($\tau$ = 0.9647), despite substantial accuracy degradation (~12\% mean drop).

\#\#\#\# Uniform Degradation Mechanism

The high ranking correlation is explained by uniform accuracy degradation. All models---regardless of architecture, parameter count, or publication year---experience similar percentage drops when evaluated on ImageNet-V2. This uniformity preserves ordinal relationships: if Model A outperforms Model B on ImageNet by 2 percentage points, it typically still outperforms Model B on V2 by approximately 2 percentage points.

\subsubsection{Limitations}

\textbf{Only One Benchmark Pair Tested.} The analysis is limited to ImageNet $\rightarrow$ ImageNet-V2. Other distribution shifts (ObjectNet, ImageNet-Sketch, ImageNet-R) may show different patterns. ImageNet-V2 was constructed to replicate the original data collection methodology, making it a relatively near distribution shift. The observed ranking stability may partially reflect this design choice rather than inherent model robustness.

\textbf{Sample May Not Represent Full Population.} The analysis covers 96 models with reported results on both benchmarks. Selection bias may exist: models expected to generalize well may be overrepresented in V2 reporting.

\textbf{Original Hypothesis Refuted.} The hypothesis that $\tau$ < 0.90 was not supported; $\tau$ = 0.96 was observed instead. While this is a null result with respect to the original prediction, it establishes that ranking instability is not a significant concern for ImageNet-V2.

\subsubsection{Broader Implications}

Validating benchmark-based model selection reduces the risk of practitioners choosing suboptimal models due to misleading benchmarks. For distribution shifts similar to ImageNet-V2, model selection decisions based on ImageNet rankings remain valid.

